% !TeX program = xelatex
% UTF-8 transcription of all five source pages, in handwritten page order.
% Original wording, numbering, abbreviations and highlighted passages retained.
% pic01.jpg = 20261019_081554281.jpg
% pic02.jpg = 20261019_081559342.jpg
% pic03.jpg = 20261019_081612704.jpg
% pic04.jpg = 20261019_081617739.jpg
% pic05.jpg = 20261019_081627378.jpg
% Compile directly with XeLaTeX, or input into a document using XePersian.
\ifdefined\settextfont
  \def\NAPRStandalone{0}
\else
  \def\NAPRStandalone{1}
  \documentclass[11pt,a4paper]{article}
  \usepackage[margin=1.8cm]{geometry}
  \usepackage{amsmath,amssymb,xcolor,enumitem}
  \usepackage{xepersian}
  \IfFontExistsTF{Amiri}{\settextfont{Amiri}}{\settextfont{FreeSerif}}
  \setlatintextfont{FreeSerif}
\fi
\providecommand{\en}[1]{\lr{#1}}
\newcommand{\NAPRhl}[1]{\ifdefined\colorbox\colorbox{orange!40}{\strut #1}\else\textbf{#1}\fi}
\ifdefined\setlist\setlist[itemize]{itemsep=0.15em,topsep=0.3em}\fi
\newcommand{\NAPRhead}[1]{\clearpage
  \noindent\en{Note book}\hfill\en{Subject 13}\par
  \noindent\en{Veor:}\hfill\en{Month:}\hfill\en{Date:}\par
  \medskip\noindent\en{Page #1}\par\medskip}
% The printed notebook label reads "Veor:" rather than "Year:".
\newcommand{\NAPRfoot}[1]{\par\medskip\noindent
  \en{BAHARE DANESH ASL}\hfill #1\par}
\ifnum\NAPRStandalone=1 \begin{document}\fi

\NAPRhead{1}
\begin{latin}
\section*{Normal Abdomen and Pelvic on conventional radiographs:}
\subsection*{1. What to look at}
\begin{itemize}
\item Overall gas pattern
\item Extraluminal air
\item Abnormal abdominal calcification
\item Soft tissue masses
\end{itemize}
\subsection*{2. Normal bowel gas pattern}
\begin{itemize}
\item \textbf{All} gas in the bowel comes from swallowed air.
\item Loops of bowel that contain a sufficient amount of air to be distended.
  \begin{itemize}
  \item \NAPRhl{Distention of bowel is normal}
  \end{itemize}
\item Loops of bowel that are filled beyond their normal size are dilated.
  \begin{itemize}
  \item \NAPRhl{Dilation of the bowel is abnormal}
  \end{itemize}
\item There is always air in the stomach unless:
  \begin{itemize}
  \item The patient has recently vomited
  \item There is a NG-Tube attached to suction
  \end{itemize}
\item There is usually a small amount of air in \NAPRhl{2 or 3 loops of small bowel}.
  \begin{itemize}
  \item Normal diameter of small bowel $<2.5\,\mathrm{cm}$
  \end{itemize}
\item There is always air in the rectum and sigmoid.
\item Varying amount of gas in the remainder of the colon.
\item Stool: Multiple, small bubbles of gas within a semisolid appearing soft-tissue mass.
\end{itemize}
\end{latin}
\noindent ۳. در گرافی \en{prone} کولون صعودی و نزولی بالاتر قرار دارند و هوا در آن‌ها تجمع پیدا می‌کند.
\NAPRfoot{۱}

\NAPRhead{2}
\begin{latin}
\subsection*{4. Normal fluid level}
\begin{itemize}
\item There is always air-fluid level in the stomach.
\item \NAPRhl{2 or 3 air-fluid} levels in the small bowel may be seen.
\item There is usually no or very few air fluid levels in the colon.
\end{itemize}
\subsection*{5. DDX for many air-fluid levels in the colon}
\begin{itemize}
\item Recent enema
\item Medications with
  \begin{itemize}
  \item Anticholinergic effect
  \item Antiperistaltic effect
  \end{itemize}
\end{itemize}
\end{latin}
\subsection*{۶. افتراق رودهٔ بزرگ از رودهٔ کوچک}
\begin{latin}
\begin{itemize}
\item Large bowel is peripherally placed
\item Haustral markings usually do not extend completely across the large bowel.
\item Haustral markings are spaced more widely apart than valvulae conniventes of the small bowel
\item \NAPRhl{Small bowel max dilated diameter: 5 cm}
\end{itemize}
\subsection*{7. Acute abdomen series}
\begin{enumerate}[label=\alph*)]
\item Supine view (Scout Film)
\item Prone or lateral rectum view
\item Upright or Left-lateral decubitus view
\item Chest: upright or supine view
\end{enumerate}
\end{latin}
\NAPRfoot{۲}

\NAPRhead{3}
\begin{latin}
\subsection*{8. Supine view (Scout Film)}
\begin{itemize}
\item Overall gas pattern
\item Calcifications
\item Masses
\end{itemize}
\subsection*{9. Prone abdomen view}
\begin{itemize}
\item Gas in the rectosigmoid
\end{itemize}
\subsection*{10. Upright abdomen view}
\begin{itemize}
\item \NAPRhl{Free air}
\item \NAPRhl{Air-Fluid levels} in the bowel
\end{itemize}
\subsection*{11. Upright chest}
\begin{itemize}
\item Free air
\item Pneumonia
\item Pleural Effusion
\end{itemize}
\subsection*{12. Causes of PE in the abdomen}
\begin{itemize}
\item Pancreatitis (left pleural effusion)
\item Ovarian tumors (right sided or bilateral pleural effusion)
\item Abscess beneath the right hemidiaphragm (right pleural effusion)
\end{itemize}
\subsection*{13. Non-pathologic calcifications of the abdomen}
\begin{itemize}
\item Phleboliths
  \begin{itemize}
  \item Small, rounded calcifications.
  \item Calcified venous thrombi
  \item Occures with aging.
    \begin{itemize}
    \item In the pelvic veins of \textbf{woman}
    \end{itemize}
  \item They have a lucent center
  \end{itemize}
\item Calcification of the rib cartilages
  \begin{itemize}
  \item Occures with aging
  \item Amorphous, speckled appearance
  \item DDX: Renal or Billiary calculi
  \item More common in \textbf{Female}. Rare under the age of \textbf{35}
  \end{itemize}
\end{itemize}
\end{latin}
\NAPRfoot{۳}

\NAPRhead{4}
\begin{latin}
\subsection*{14. Organomegaly on conventional radiograph}
\subsubsection*{a) Liver}
\begin{itemize}
\item Normal
  \begin{itemize}
  \item The liver normally displaces all bowel gas from the RUQ.
  \item A \NAPRhl{tongue-like projection} of the right lobe may extend to the iliac crest.
    \begin{itemize}
    \item Specially in women
    \item Called \NAPRhl{Riedel lobe} and is normal
    \end{itemize}
  \end{itemize}
\item Enlarged liver
  \begin{itemize}
  \item \NAPRhl{use CT, MRI, US}
  \item Displacement of all bowel from RUQ to the iliac crest and across the midline
  \end{itemize}
\end{itemize}
\subsubsection*{b) Spleen}
\begin{itemize}
\item Normal
  \begin{itemize}
  \item in adults is \NAPRhl{about 12 cm}.
  \item Does not project below the 12th posterior rib.
  \item The spleen is about as large as the left Kidney.
  \end{itemize}
\item Enlarged spleen
  \begin{itemize}
  \item If the spleen projects below the \NAPRhl{12th posterior rib}
  \item If it displaces the stomach bubble toward or across the midline.
  \end{itemize}
\end{itemize}
\subsubsection*{c) Kidney}
\begin{itemize}
\item Normal
  \begin{itemize}
  \item The Kidney length is the height of \NAPRhl{4 lumbar} \NAPRhl{vertebral bodies} or about \NAPRhl{10--14 cm}.
  \end{itemize}
\end{itemize}
\end{latin}
\NAPRfoot{۴}

\NAPRhead{5}
\begin{latin}
\subsection*{15. Enlarged bladder}
Enlarged bladder is usually recognized by displacement of bowel upward and out of the pelvis by a mass of soft tissue density
\end{latin}
\bigskip
\noindent ۱۶. در مردان فضای بالای مثانه را کولون سیگموئید پر می‌کند ولی در زنان می‌تواند رحم یا کولون سیگموئید آنجا باشد.
\par\bigskip
\noindent ۱۷. \NAPRhl{اندازهٔ نرمال رحم \en{$4\times6\times8$} سانتی‌متر است.}
\par\bigskip
\noindent ۱۸. بهترین تصویربرداری برای بررسی بزرگ شدن تخمدان و رحم \en{$\longleftarrow$ US}
\NAPRfoot{۵}

\ifnum\NAPRStandalone=1
  \def\NAPRfinish{\end{document}}
\else
  \let\NAPRfinish\relax
\fi
\NAPRfinish
